% Ejercicio 1.
Para crear los ficheros de sonido en la página \url{https://speechgen.io/} hay que seleccionar el idioma y la voz (arriba a la izquierda). A continuación hay que escribir el texto en el recuadro. Observar Figura \ref{fig:img1}.\\ 

\begin{figure}[htbp]
    \centering
    \includegraphics[width=0.9\linewidth]{images/1_1convertir_nombre1.png}
    \caption{Convertir texto a sonido}
    \label{fig:img1}
\end{figure}

Abajo a la izquierda hay un botón de ``SETTINGS'' (Figura \ref{fig:img2}) que sirve para cambiar la configuración del archivo.

\begin{figure}[H]
    \centering
    \includegraphics[width=0.9\linewidth]{images/1_2convertir_ajustes.png}
    \caption{Convertir texto a sonido}
    \label{fig:img2}
\end{figure}

Seleccionar formato ``wav''.\\

Por último clikar en ``CONVERT TO SPEECH''. Aparecerá abajo el audio, y podemos descargarlo pulsando el botón ``DOWNLOAD''. Observar Figura \ref{fig:img3}.

\begin{figure}[H]
    \centering
    \includegraphics[width=0.9\linewidth]{images/1_3descargar.png}
    \caption{Descargar sonido}
    \label{fig:img3}
\end{figure}

Para crear el fichero de sonido con los apellidos hay que realizar el mismo procedimiento. Observar Figura \ref{fig:img4}.
\begin{figure}[H]
    \centering
    \includegraphics[width=0.9\linewidth]{images/1_4apellidos.png}
    \caption{Fichero de sonido con apellido}
    \label{fig:img4}
\end{figure}
